% !TEX encoding = UTF-8 Unicode
\documentclass[a4paper]{article}

\usepackage{color}
\usepackage{url}
\def\UrlBreaks{\do\/\do-}
\usepackage[T2A]{fontenc} % enable Cyrillic fonts
\usepackage[utf8]{inputenc} % make weird characters work
\usepackage{graphicx}
\usepackage{makecell}
\usepackage{pgfplots}

\usepackage[english,serbian]{babel}
%\usepackage[english,serbianc]{babel} %ukljuciti babel sa ovim opcijama, umesto gornjim, ukoliko se koristi cirilica

\usepackage[unicode]{hyperref}
\hypersetup{colorlinks,citecolor=green,filecolor=green,linkcolor=blue,urlcolor=blue}

\usepackage{listings}

%\newtheorem{primer}{Пример}[section] %ćirilični primer
\newtheorem{primer}{Primer}[section]

\definecolor{mygreen}{rgb}{0,0.6,0}
\definecolor{mygray}{rgb}{0.5,0.5,0.5}
\definecolor{mymauve}{rgb}{0.58,0,0.82}

\lstset{ 
  backgroundcolor=\color{white},   % choose the background color; you must add \usepackage{color} or \usepackage{xcolor}; should come as last argument
  basicstyle=\scriptsize\ttfamily,        % the size of the fonts that are used for the code
  breakatwhitespace=false,         % sets if automatic breaks should only happen at whitespace
  breaklines=true,                 % sets automatic line breaking
  captionpos=b,                    % sets the caption-position to bottom
  commentstyle=\color{mygreen},    % comment style
  deletekeywords={...},            % if you want to delete keywords from the given language
  escapeinside={\%*}{*)},          % if you want to add LaTeX within your code
  extendedchars=true,              % lets you use non-ASCII characters; for 8-bits encodings only, does not work with UTF-8
  firstnumber=1000,                % start line enumeration with line 1000
  frame=single,	                   % adds a frame around the code
  keepspaces=true,                 % keeps spaces in text, useful for keeping indentation of code (possibly needs columns=flexible)
  keywordstyle=\color{blue},       % keyword style
  language=Python,                 % the language of the code
  morekeywords={*,...},            % if you want to add more keywords to the set
  numbers=left,                    % where to put the line-numbers; possible values are (none, left, right)
  numbersep=5pt,                   % how far the line-numbers are from the code
  numberstyle=\tiny\color{mygray}, % the style that is used for the line-numbers
  rulecolor=\color{black},         % if not set, the frame-color may be changed on line-breaks within not-black text (e.g. comments (green here))
  showspaces=false,                % show spaces everywhere adding particular underscores; it overrides 'showstringspaces'
  showstringspaces=false,          % underline spaces within strings only
  showtabs=false,                  % show tabs within strings adding particular underscores
  stepnumber=2,                    % the step between two line-numbers. If it's 1, each line will be numbered
  stringstyle=\color{mymauve},     % string literal style
  tabsize=2,	                   % sets default tabsize to 2 spaces
  title=\lstname                   % show the filename of files included with \lstinputlisting; also try caption instead of title
}

\begin{document}

%\title{Desavanja u Beogradu: 8 udruzenja koja organizuju redovne IT strucne sastanke (meet up) \\ \small{Seminarski rad u okviru kursa\\Metodologija stručnog i naučnog rada\\ Matematički fakultet}}

\title{IT organizacije koje priređuju redovne stručne sastanke\\ \small{Seminarski rad u okviru kursa\\Metodologija stručnog i naučnog rada\\ Matematički fakultet}}

\author{Đorđe Vujinović, Nebojša Koturović, Igor Radojević, Bojan Stefanović\\ ai15245@alas.matf.bg.ac.rs, mi15139@alas.matf.bg.ac.rs, \\ mi18493@alas.matf.bg.ac.rs, mi15022@alas.matf.bg.ac.rs}

%vujinovic.djordje@gmail.com
%\date{1.~april 2020.}

\maketitle

\abstract{
%TODO(sažetak....)
Kroz ovaj rad čitalac će se upoznati sa nekim od najpoznatijih IT organizacija u Beogradu koje organizuju redovne susrete. Za svaku organizaciju ukratko su predstavljeni ciljevi, mesto održavanja sastanaka, razlozi održavanja, kao i istorijat same organizacije. Saznaćete za koga su okupljanja namenjena, ko su predavači i koliko učesnika ima. Za svako okupljanje navedeni su načini kako se može prijaviti u svojstvu predavača/slušaoca, kao i adrese ka sajtovima organizacija.

\textbf{Ključne reči:} \textit{organizacija, udruženje, sastanak, php, wordpress, javascript, python, agile, risk, crypto}

\tableofcontents

\newpage
\section{Uvod}

% Pokušaćemo da Vas u ovom seminarskom radu upoznamo sa največim i najbitnijim IT organizacijama koje na teritoriji Beograda organizuju redovne susrete.

Susret (eng. meetup) je okupljanje ljudi sa istim intresovanjem za neku određenu temu. Obično se sastoji od formalnog dela tj, nekog predavanja i diskusije o toj temi, zatim neke radionice, koja ne mora nužno postojati, dok je kraj ostavljen za druženje i osveženje. To je šansa za povezivanje ljudi, ostvarivanje biznis kontakata i stvaranje novih ideja. Dosta kompanija vide šansu u susretima da doprinesu svom ugledu i brendu i tako privuku nove kandidate za svoj tim.


\section{PHP Serbia}
\begin{figure}[h!]
\begin{center}
\includegraphics[scale=0.25]{php_gray.jpg}
\end{center}
\caption{PHP Srbija logo}
\label{fig:pande}
\end{figure}
PHP je popularan skriptni programski jezik za izradu dinamičkih web stranica. Prema mnogim listama, PHP spada u top 10 najpopularnijih programskih jezika \cite{phpMostPopular}. U Srbiji takođe spada među najpopularnije jezike, što opravdava veliku zajednicu i višegodišnje okupljanje \cite{phpSerbiaPopularity}.

Prvi PHP Serbia Meetup organizovan je u julu 2015. godine u KC Gradu, i tadašnji plan je bio da se okupljanja održavaju svakog meseca. Do sada je organizovano 38 okupljanja u Beogradu, ne računajući radionice i panele. Predavači su uglavnom PHP programeri i sistem administratori, a teme koje se obrađuju su vezane za sam jezik, alate, framework-e, kao i iskustva i probleme sa kojima programeri mogu da se susretnu. Neke teme i predavače možete videti u Tabeli \ref{tab:tabelaPHP}.  Na kraju svakog okupljanja slede pitanja posetilaca ili panel diskusija na neku temu.\\
Broj posetilaca na okupljanjima varira, a na nekima je prešao broj od 160. Posećenost na okupljanjima možete videti na grafikonu ispod \cite{phpEvents}. 

Iza samog okupljanja stoji organizacija PHP Srbija koja je zaslužna i za PHP Srbija konferenciju koja je dovele neke od najboljih svetskih predavača, kao i kreatora programskog jezika PHP Rasmusa Lerdorfa \cite{phpRasmusLerdorf}.

Forma za prijavljivanje u svojstvu predavača se može naći na zvaničnom sajtu PHP Srbije \url{https://phpsrbija.rs/}. Prijavljivanje u svojstvu slušaoca je besplatno, samo je potrebno prijaviti se na \url{https://meetup.com/PHPSrbija}.

Iz ovog okupljanja stvorilo se još nekoliko samostalnih kao na primer Laravel Serbia Meetup. Na PHP Srbija youtube kanalu se mogu naći predavanja sa okupljanja i konferencija, tako da svako ko nije bio u mogućnosti da dođe može da odgleda. Adresa youtube kanala je: \url{https://www.youtube.com/user/PHPSrbijaVideo/}.


\begin{table}[h!]
\begin{center}
\caption{Neke od tema sa okupljanja.}
\begin{tabular}{|l|l|l|} \hline
\thead{Tema} & Predavač&\#. okupljanje\\ \hline
FigDice template engine&NIkola Posa&1\\ \hline
JMeter i Performance Testing&Nebojša Videnov&6\\ \hline
\makecell[l]{Beginner talk: The road to\\become a junior developer}&\makecell[l]{Vladimir Živadinović, \\ Operations manager}&17\\ \hline
\makecell[l]{PHP Aplikacije u \\produkcionom okruženju}&\makecell[l]{Nikola Krgović,\\Sistem administrator}&28\\ \hline
\makecell[l]{Razvoj efikasnih API servisa - \\Laravel i GraphQL}&\makecell[l]{Peđa Jevtić,\\Full stack developer}&37\\ \hline
\end{tabular}
\label{tab:tabelaPHP}
\end{center}
\end{table}

\begin{tikzpicture}
\begin{axis}[
    title={Posećenost PHP Serbia Meetup-a},
    xlabel={Redni broj okupljanja},
    ylabel={Broj posetilaca},
    xmin=0, xmax=40,
    ymin=0, ymax=180,
    xtick={0,5,10,15,20,25,30,35,40},
    ytick={0,20,40,60,80,100,120,140,160},
    legend pos=north west,
    ymajorgrids=true
]

\addplot[
    color=blue,
    mark=square,
    ]
    coordinates {
    (1,55)(2,83)(3,86)(4,113)(5,110)(6,89)(7,45)(8,72)(9,66)(10,82)(11,64)(12,94)(13,82)(14,35)(15,73)(16,137)(17,165)(18,137)(19,133)(20,159)(21,116)(22,96)(23,137)(24,56)(25,139)(26,49)(27,91)(28,63)(29,95)(30,80)(31,62)(32,35)(33,88)(34,59)(35,41)(36,50)(37,88)(38,42)
    };
    
\end{axis}
\label{fig:M1}
\end{tikzpicture}

\section{WordPress Serbia}
\begin{figure}[h!]
\begin{center}
\includegraphics[scale=0.5]{wp.jpg}
\end{center}
\caption{WordPress Serbia logo}
\label{fig:pande}
\end{figure}
WordPress predstavlja jedan od najpopularnijih sistema za upravljanje sadržajem (eng.~{\em Content Management System - CMS}). Zbog lakoće podešavanja, brzog rasta platforme, mogućnosti kreiranja dodataka i tema kao i korišćenje velikog broja besplatnih iz njihovog marketa, ubrzo je pridobio pažnju programera i entuzijasta u Srbiji. WordPress Serbia je organizacija koja za cilj ima da ih sve okupi na jednom mestu. 

Prvi WordPress Serbia Meetup održan je u aprilu 2013. godine. Sledeće godine održavaju se ukupno 4 okupljanja, od čega su prva tri održana u Mikser House-u, dok je poslednji održan u Domu omladine u tribinskoj sali. Tada počinje velika zainteresovanost, tako da se prelazi u veliku salu Doma omladine. Ubrzo nastaje potreba širenja zajednice, tako da okupljanja počinju i u Nišu, Inđiji, Šapcu, Boru, Vranju, Kruševcu, Subotici, Novom Sadu, Zrenjaninu itd. Pored okupljanja kreću i sa radionicama za početnike koje su namenjene svima koji žele da nauče više o WordPress-u.

Teme koje se obrađuju na okupljanjima su prilagođene početnicima, programerima, kao i ostatku IT zajednice. Neke od tema sa okupljanja se nalaze u Tabeli \ref{tab:tabelaWordpress}. \\ \\
\begin{table}[h!]
\begin{center}
\caption{Neke od tema sa okupljanja.}
\begin{tabular}{|l|l|l|} \hline
\thead{Tema} & Predavač& \#. okupljanje\\ \hline
WordPress security & \makecell[l]{Predrag Cujanovic - \\CEO, CyberTec Security}&2\\ \hline
\makecell[l]{Scaling WordPress \\with Amazon cloud} &\makecell[l]{Miljenko Rebernisak,\\Devana Technologies}&9\\ \hline %TODO(nkoturovic: Ovde je bio neki cudan karakter zbog kog nije htelo da se kompajlira!)
\makecell[l]{Optimizacija veb-sajta za \\pretraživače - Pravila za \\dobru optimizaciju} &Stevica Gološin&20\\ \hline
\end{tabular}
\label{tab:tabelaWordpress}
\end{center}
\end{table}


Predavači su najčešće programeri, firme koje se bave izradom tema i dodataka za WordPress, dizajneri, kao i sistem administratori. Za predavača može da se prijavi bilo ko ko želi da podeli svoje iskustvo i znanje. 

Prijavljivanje za predavača je putem google forme na adresi \url{https://docs.google.com/forms/d/1kHF5Vi-35FkmAkRzdsHiQOoKrIPGOjKgVRE8mAgjA-g}. Prijavljivanje u svojstvu slušaoca je besplatno, potrebno je samo da se na adresi \url{https://www.meetup.com/WordPress-Serbia/} prijavi na željeni događaj.

Pored Meetup-a, WordPress Serbia organizuje i WordCamp konferencije na kojima gostuju svetski poznati predavači, kao i posetioci iz čitave Evrope i sveta \cite{wpWordCamp}.

\section{Python Belgrade}
\label{sec:pybgd}

\begin{figure}[h!]
  \centering
  \includegraphics[width=0.3\textwidth]{pybgd.png}
  \caption{Python Belgrade logo}
\end{figure}

Python Belgrade je nevladina organizacija koja organizuje okupljanja za sve Python entuzijaste. Ova okupljanja namenjena su svim
zainteresovanima za jezik Python, pre svega programerima, ali i svim onima koji žele da upotpune svoje znanje ovog programskog jezika i njegovih mogućnosti.

Jezik Python je nesumnjivo jedan od najpopularnijih programskih jezika danas. Sudeći po istraživanju kompanije TIOBE, jezik Python ima sigurno mesto u top 10 najpretraživanijih jezika na poznatim web pretraživačima. \cite{pythonPopular}.

Kako i samo ime organizacije kaže, sastanci se organizuju na teritoriji grada Beograda. Ova organizacija nema svoje prostorije, a lokacija, kao i ostale informacije
za sledeći susret dostupne su na oficijalnom vebsajtu organizacije (\url{https://pythonbelgrade.com/}). Ova oraganizacija je svakako jedna od značajnijih IT organizacija koje organizuju okupljanja na teritoriji Beograda. Osnovana je 2015. godine i od tada se predavanja održavaju svakih par meseci bez striktnog međuintervala. 

Na čelu ove organizacije, ispred grupe mladih ljudi, je njen domaćin i osnivač \textit{Bojan Jovanović}. Bojan je mladi entuzijasta koji je i sam imao dosta radnog iskustva u ovom jeziku nakon završenih studija na
Elektrotehničkom fakultetu. Pored organizovanja Python Belgrade okupljanja, on je i jedan od osnivača PyCon Balkan konferencije (\url{https://pyconbalkan.com/}). Ova konferencija važi za jednu od najpoznatinjih konferencija vezanih za programiranje na ovim prostorima.

Predavači na Python Belgrade okupljanjima su najčešće ljudi koji su zaposleni ili imaju neke veze sa IT sektorom, ali neretko i ljudi sa univerziteta. \cite{pybgdSpeakers} Iako se predavanja održavaju na srpskom jeziku, dešavalo se da zbog prisustva stranih slušalaca u duhu solidarnosti bude napravljen izuzetak i predavanje se održi na engleskom jeziku. Na okupljanjima se nesumnjivo oseća mladalački duh i želja za napredovanjem. Uvek vlada prijateljska atmosfera, nađu se tu grickalice i osveženje, pa nije ni čudo što se okupljeni između sebe zovu ``Pytosi''.

Putem vebsajta \url{https://pythonbelgrade.com/} se može prijaviti za prisustvo nekom od okupljanja, a obaveštenja se objavljuju i na fejsbuk grupi \href{https://www.facebook.com/python.belgrade/}{@python.belgrade} i \href{https://pythonbelgrade.slack.com/join/shared_invite/enQtNTM0OTI4NzY3NDI2LTk2OThkNjQ2YmM5NWNkNTBjMjU3NDY0NjcyNDFiYmZmNjQ4MDE3NjczYWM0NTk1YWVlYzlhNGIwYjU3NjA3Y2Y}{slack kanalu}. Svako ko želi da se prijavi u svojstvu predavača potrebno je da popuni odgovarajuću web formu za prijavu na \href{https://pythonbelgrade.com/}{vebsajtu} organizacije (\url{https://pythonbelgrade.com/#contact}). Predavanja se mogu pratiti i putem video linka uživo, gde se takođe mogu postavljati pitanja. Neka od već održanih predavanja dostupna su na YouTube kanalu ove organizacije (\url{https://www.youtube.com/channel/UCH79IYl8rv2f2BrYs5owB2A}).


\section{RISK}
\label{sec:riskmatf}

\begin{figure}[h!]
  \centering
  \includegraphics[width=0.4\textwidth]{riskmatf.png}
  \caption{RISK logo}
\end{figure}

RISK predstavlja organizaciju koja za cilj ima okupljanje i edukaciju studenata zainteresovanih za oblast računarstva i informatike. \cite{aboutRisk}
Ova organizacija ima svoju web stranu, međutim informacije su dostupne i na facebook grupi, gde se unapred najavljuju okupljanja \url{https://www.facebook.com/groups/1192849077458638/}. Predavanju može prisustvovati bilo ko, mada je poželjno da se prijavi putem web forme koja je dostupna pre svakog okupljanja. Najveći broj slušalaca su
studenti Matematičkog, ali tu su i kolege sa ETF-a, FON-a i ostalih srodnih fakulteta.

Udruženje je osnovano u decembru 2016. godine, kada je i održano prvo okupljanje. Okupljanja se organizuju u zgradi ``Matematičkog fakulteta'', a do danas ih je održano preko 20. Najčešće se održavaju vikendom, kada postoje mogućnosti za to. Predavač može biti bilo ko zainteresovan i voljan da podeli svoje programersko umeće. Na čelu ove studentske organizacije je Nemanja Mićović, koji je i sam bio student Matematičkog fakulteta, a sada je asistent i pohađa doktorske studije. Pored Nemanje, u organizaciji učestvuju i studenti sa osnovnih i master studija (\url{http://risk.matf.bg.ac.rs/organizacija.html}).

Na predavanjima se obrađuju razne teme vezane za programiranje. Neretko su to radni okviri za razvoj aplikacija, programerski alati i druge tehnologije.
Pored ovih okupljanja, organizacija može da se pohvali i inspirativnim blog postovima koji su izuzetno zanimljivi. (\url{http://risk.matf.bg.ac.rs/blog.html}).

Svoj doprinos ovoj organizaciji možete ostaviti u vidu saveta, kritika i predloga popunjavanjem odgovarajuće forme na sajtu. Ukoliko ste voljni, možete se uključiti
i u neki od projekata ove organizacije čiji su izvorni kodovi dostupni na github-u \url{https://github.com/riskmatf/}.

\begin{table}[h!]
\caption{Neke od tema sa okupljanja}
\begin{center}
\begin{tabular}{|l|l|l|} \hline
\thead{Tema} & Predavač& Datum\\ \hline
\makecell[l]{Wifi Hacking}&\makecell[l]{Hacklab Beograd}&23.2.2019.\\ \hline
\makecell[l]{Uvod u duboko učenje kroz\\PyTorch}&\makecell[l]{Nemanja Mićović}&3.11.2019.\\ \hline
\makecell[l]{Moderno Android Programiranje}&\makecell[l]{Aleksandar Stefanović}&10.3.2019.\\ \hline
\makecell[l]{Uvod u razvoj veb aplikacija\\kroz Flask radni okvir.}&\makecell[l]{Stevan Nestorović}&18.3.2018.\\ \hline
\makecell[l]{Uvod u BASH skripting}&\makecell[l]{Peđa Trifunov}&2.12.2017.\\ \hline
\makecell[l]{Uvod u blockhain tehnologije\\i decentralizaciju.}&\makecell[l]{DECENTER,\\origintrail}&24.12.2017.\\ \hline
\makecell[l]{Git i Github}&\makecell[l]{Marko Jeremić}&12.12.2016.\\ \hline

\end{tabular}
\label{tab:tabelaRISK}
\end{center}
\end{table}



\section{Belgrade Crypto Community}
\label{sec:bgdcs}

\begin{figure}[h]
  \centering
  \includegraphics[width=0.3\textwidth]{bcc_logo.png}
  \caption{Belgrade Crypto Community logo}
  \label{fig:bgdcclogo}
\end{figure}

Termini  kriptovaluta, blokčejn, majner, digitalni novčanik i bitkoin su od svog nastanka zaintrigirali i dobili pažnju ne samo IT scene, već i javnosti. To je bilo lakši deo, dok tačno definisanje pojmova, koncepta i relacija među njima baš i ne. Belgrade Crypto Community je organizacija sa ciljem da ovu tematiku popularizuje i prenese svoje znanje na što veći broj ljudi.

Belgrade Crypto je domaća kripto zajednica koja povezuje više od 1200 članova, od korisnika i novajlija do iskusnih programera blokčejna. Njihov logo možete videti na slici \ref{fig:bgdcclogo}. Prostorije im se nalaze u Balkanskoj 2, gde se održavaju okupljanja, dobro su opremljene i imaju kapaciteta za veći broj učesnika. Informacije o samoj zajednici kao i buduće susrete možete naći na zvaničnom veb-sajtu \url{https://belgradecrypto.com/}. Na sajtu postoji blog koji se bavi zanimljivim i aktuelnim temama.

Prvi Belgrade Crypto Meetup održan je u februaru 2017. godine. Od tad su promenili više lokacija i formata samog okupljanja, trenutno se održavaju svakog drugog četvrtka u mesecu od 18 do 20 časova. U planu je organizacija i u drugim gradovima, a prvi takav je bio u Novom Sadu. Susreti se obično realizuju na jedan od dva načina: 1. opušteno druženje, pričanje o aktulenostima i novitetima; 2.predavanje eksperata. Bilo je i nekoliko radionica, na kojima su učesnici mogli  da se upoznaju sa lokalnom berzom Xcalibra, globalnom platformom za trgovinu Safex i automatom za kupovinom kritovaluta HashBringer. 

Forme za prijavu u svojstvo predavaca ili slusaoca je moguče naći na stranici \url{https://www.meetup.com/Belgrade-Crypto-Community/}. Belgrade Crypto Community ima stranice na društvenim mrežama \url{https://www.facebook.com/BelgradeCryptoCommunity/}, \url{https://twitter.com/Belgrade_Crypto}, \url{https://www.instagram.com/belgrade_crypto/}, \url{https://t.me/belgrade_crypto}, preko kojih informiše o svojim aktivnostima. Pojedine susrete i video tutoriale možete naći na You-Tube stranici \url{https://www.youtube.com/channel/UCCv76WWtH9LppS65OCWpI4A/featured}.


\section{Agile Serbia}
\label{sec:agsrb}

\begin{figure}[h]
  \centering
  \includegraphics[width=0.3\textwidth]{agile_srb.png}
  \caption{Agile Serbia logo}
  \label{fig:agslogo}
\end{figure}

Agilni razvoj softvera je skup okvira i metoda razvoja softvera na principima iterativnog i inkrementalnog razvoja, gde zahteve korisnika u rešenja implementiraju samoorganizovani razvojni timovi.Agile Manifesto je nastao u februaru 2001. godine, u gradu Juta, kad su se sastali 17 programera da razgovaraju o lakšim metodama razvoja. Agilne poslovanja rade brže, bolje, sa većom vrednošću izlaznog proizvoda i sa većim procentom uspešno završenih projekata.

Agile Serbia je nastala kao edukativni centar kompanije Puzzle Software 2012 godine. Od tad im je glavni cilj agilni razvoj softvera i njegovo unapređenje.Osnivači imaji bogato iskustvo sa tradicionalnim metodama razvoja softvera i borili se sa izazovima i manama tih metoda. Novi pristup je bio Skram čije korisčenje je delovalo sasvim logično tako da je novi čilj pomoći drugima sa istim problemima i proširiti vec postojeće znanje. U istu svrhu je nastao i blog (\url{http://www.agile-serbia.rs/blog/}). Logo Agile Serbia možete videti na slici \ref{fig:agslogo}.

Sama organizacija se nalazi na Banjici, a susrete generalno održavaju u halama prestoničkih hotela, startit ili ict centru.Određen broj susreta je bio u Novom Sadu i Nisu. Od 2012 okupljanja se organizuju fleksibilno u skladu sa agile duhom, nekad jednom u mesecu, nekad jednom u tri meseca, stime sto mesec maj obično bude mesec agilnog razvoja i tad ima dosta dešavanja. Predavači su ljudi sa dugogodišnjim iskustvom iz zemlje, a i iz inostranstva. Posebno su ponosni na činjenicu da je nekoliko ko-kreatora agilnog manifesta održalo predavanje na njihovim skupovima. Neke teme predavače i broj učesnika možete videti u Tabeli \ref{tab:tabelaAgS}.

Pored susreta, Agile Serbia pravi jednom godisnje konferenciju Agile Serbia Conference. Četvrto izdanje se održalo sredinom juna 2019 u prostorijama Madlenianuma i zvanični moto je bio "Agile, In \& Beyond IT". Bilo je preko 550 učesnika, 14 predavača na 3 bine i preko 45 zvaničnih partnera. Među predavačima je bio Jeff Sutherland, ko-kreator agilnog manifesta i skrama, Sander Hoogendoorn, Riina Hellstrom, pionir agila u hr, Milos Zeković, Ivan Avdić, Predrag Rajković i mnogi drugi.\\

\begin{table}[h!]
\caption{Neke od tema sa okupljanja}
\begin{center}
\begin{tabular}{|l|l|l|l|} \hline
\thead{Tema} & Predavač & Datum & Br. učesnika\\ \hline
\makecell[l]{Product Owner as a Product \\Mana	ger and Product Developer}&\makecell[l]{Predrag Rajković}&17.10.2019. & 60\\ \hline
\makecell[l]{The 4 Dynamics of Agility \\that Every Scrum Master and \\Leader Should Know}&\makecell[l]{Miljan Bajić}&10.09.2019. & 60\\ \hline
\makecell[l]{Uvod u Agile i Scrum}&\makecell[l]{Miloš Zeković}&28.03.2019. & 60\\ \hline
\makecell[l]{Da li Design Sprint ima smisla\\u korporaciji?}&\makecell[l]{Predrag Rajković i\\Zoran Vujkov}&30.10.2018. & 21\\ \hline
\makecell[l]{Kako biti lider u napetoj radnoj\\sredini?}&\makecell[l]{Olaf Lewitz}&25.10.2018. & 70\\ \hline
\makecell[l]{Gasi teoriju! Pali praksu!}&\makecell[l]{Davor Čengija i\\Ivan Krnić}&05.12.2017. & 93\\ \hline
\makecell[l]{Da li ste dovoljno agilni?}&\makecell[l]{Zoran Vujkov}&01.10.2016. & 46\\ \hline
\end{tabular}
\label{tab:tabelaAgS}
\end{center}
\end{table}

Najnovije informacije možete naći na zvaničnom sajtu (\url{http://www.agile-serbia.rs/}) ili na društvenim mrežama (\url{https://www.facebook.com/AgileSerbia/}), (\url{https://twitter.com/PuzzleSoft}), (\url{https://www.youtube.com/channel/UC01ZnqLRvcxG9CH3Ssjswkw/featured}), kače snimke čelih predavanja, (\url{https://www.linkedin.com/company/13055153/}). 


\section{JS Belgrade}
\label{sec:JSBM}

\begin{figure}[h!]
\begin{center}
\includegraphics[scale=0.25]{JS_logo.png}
\end{center}
\caption{JS Belgrade Meetup logo}
\label{fig:JS_slika}
\end{figure}

JavaScript je skript programski jezik koji je inicijalno napravljen za web programiranje na klijentskoj strani. JavaScript se interpretira i danas se moze izvrsavati i van brauzera, moze se izvsavati i na serveru, zapravo moze se izvrsavati na bilo kojem uredjaju koji ima specijalan program nazvan JavaScript engine \cite{aboutJS}. 

JS Belgrade Meetup ciji je logo prikazan na slici \ref{fig:JS_slika} je JavaScript korisnicka grupa organizovana 2014 godine od strane Slobodana Stojanovica, Bogdana Gavrilovica i Alexandera Simovica sa ciljem da povezu ljude zainteresovane za JavaScript kroz mesecna okupljanja bilo da ga oni svakodnevno koriste ili su početnici koji samo žele da saznaju nesto više o JavaScript-u \cite{JS_Belgrade_meetup_pocetna}. 

Prvo JS Blegrade okupljanje bilo je 27.12.2014 u Beogradu u prostorijama Cloud Horizona na kojoj je prisustvovalo 15 osoba i bilo je reci na temu "AngularJS - juce, danas i sutra". Do sada je organizovano 36 okupljanja u Beogradu, ne računajući radionice i druzenja u Koffein-u. Od tih 36 okupljanja, koliko je okupljanja odrzano u kojoj godini mozete videti na tabeli \ref{tab:JStabela1}, a broj prisutnih na svakom od tih okupljanja mozete videti na slici \ref{fig:JSgrafik}, a detalje o 7 najposecenijih okupljanja mozete videti na tabeli \ref{tab:JStabela2}. Okupljanja su se organizovala na razlicitim lokacijama, od kojih je vecina bila u ICT HUB prostorijama i u Cloud Horizon prostorijama, ali i na drugim lokacijama poput Matematicki fakultet, Crowne Plaza Belgrade, Startit Centar itd. Informacije o proslim i buducim okupljanjima mozete naci na linku: \url{https://www.meetup.com/JS-Belgrade-Meetup/events/}. Svako ko hoce moze da prisustvuje okupljanjima kao slusalac samo je potrebno na vreme se prijaviti jer su mesta uglavnom ogranicena. Takodje svako ko zeli da prica na okupljanju zahtev za razgovor moze podneti na linku: \url{www.jsbelgrade.org/teme/}, a upustvo za prijavu ima na linku: \url{https://github.com/JSBelgrade/cfp}. Pocev od 19.01.2016. godine svakog drugog utorka od 8 do 10 ujutru se organizuje druzenje u kaficu Koffein koji se nalazi u ulici Uskocka 8. u Beogradu sa ciljem da se okupe ljudi zainteresovani za JavaSript i upoznaju. Na tom druzenju nema predavanja ni bilo kakvog isplaniranog programa vec je to druzenje uz kafu i pricanje o raznim temama koje ne moraju biti vezane za JavaScript\cite{JS_Belgrade_meetup_events}.

Ideja JS Belgrade je da s vremena na vreme gostuje i van Beograda. Prvo takvo okupljanje je bilo u Nisu 19.03.2016 godine na kom je prisustvovalo 7 osoba. Drugo takvo okupljanje bila je radionica u Indjiji 26.03.2016. godine i to su jedina dva okupljanja van Beograda \cite{JS_Belgrade_meetup_events}.

Github repozitorijum na linku: \url{https://github.com/JSBelgrade}, mejl: \url{sbelgradeorg@gmail.com}, fejsbuk grupa na linku: \url{https://www.facebook.com/jsbelgrade/},
za diskusiju i sva dodatna pitanja o sastancima i o temama vezanim za JavaScript možete koristiti link: \url{http://slack.jsbelgrade.org}, vebsajt JS Belgrade grupe mozete naci na linku: \url{http://jsbelgrade.org/}.

\begin{table}[h!]
\begin{center}
\caption{Jos ne znam koji naslov}
\begin{tabular}{ |c|c|c|c|c|c|c|c| } 
 \hline
 godina & 2014 & 2015 & 2016 & 2017 & 2018 & 2019 & 2020 \\ 
 \hline
 \makecell[c]{br.\\semestara} & 1 & 12 & 8 & 8 & 5 & 1 & 1 \\
 \hline
\end{tabular}
\label{tab:JStabela1}
\end{center}
\end{table}

\begin{figure}
\centering
\begin{tikzpicture}
\begin{axis}[
    xlabel={Redni broj okupljanja},
    ylabel={Broj posetilaca},
    xmin=0, xmax=40,
    ymin=0, ymax=260,
    xtick={0,5,10,15,20,25,30,35,40},
    ytick={0,20,40,60,80,100,120,140,160, 180, 200, 220, 240, 260},
    legend pos=north west,
    ymajorgrids=true
]

\addplot[
    color=blue,
    mark=square,
    ]
    coordinates {
    (1,15)(2,12)(3,15)(4,19)(5,99)(7,12)(8,22)(9,110)(10,244)(11,99)(12,110)(13,25)(14,30)(15,103)(16,17)(17,29)(18,90)(19,67)(20,20)(21,30)(22,210)(23,138)(24,82)(25,96)(26,101)(27,99)(28,47)(29,70)(30,70)(31,74)(32,76)(33,80)(34,98)(35,80)(36,45)
    };
    
\end{axis}
\end{tikzpicture}
\caption{Posecenost JS Belgrade okupljanja u Beogradu(nedostaju podaci za sesto okupljenje)}
\label{fig:JSgrafik}
\end{figure}

\begin{table}[h!]
\begin{center}
\caption{Sedam najposecenijih JS Belgrade okupljanja}
\begin{tabular}{|l|l|l|l|} \hline
\thead{Tema} & \thead{Predavač} & \thead{Datum} & \thead{Br. učesnika}\\ \hline
\makecell[l]{Hoodie, npm and semantic release}&\makecell[l]{Stephan Bonnemann}&08.08.2015. & 110\\ \hline
\makecell[l]{Different Characters \\Code Golf, JavaScript, Creative coding, Hacks}&\makecell[l]{Martin Kleppe}&05.09.2015. & 244\\ \hline
\makecell[l]{Internet of Things using JavaScript for fun\\and profit}&\makecell[l]{Dejan Dimić}&28.11.2015. & 110\\ \hline
\makecell[l]{Mikroservisi (skoro) za dz: AWS + Node.JS}&\makecell[l]{Gojko Adžić}&13.02.2016. & 103\\ \hline
\makecell[l]{Theme 1: ``Async approaches: Promises\\vs Callbacks''\\Theme 2: ``Extreme performances: JS\\vs ASMjs vs PNaCl vs UWP''}&\makecell[l]{Miloš Žikić}&28.01.2017. & 210\\ \hline
\makecell[l]{Theme 1: ``Introduction to Elm''\\Theme 2: ``Developer UX, REST and the rest''}&\makecell[l]{Bojan Matić(Theme 1)\\Miloš Solujić(Theme 2)}&18.02.2017. & 138\\ \hline
\makecell[l]{``React component testing using snapshots''}&\makecell[l]{Boris Aržentar}&26.06.2017. & 101\\ \hline
\end{tabular}
\label{tab:JStabela2}
\end{center}
\end{table}

\newpage
\section{Zaključak}
\label{sec:zakljucak}

Mi, autori, iskreno se nadamo da naš rad uspešno prenosi atmosferu sa Beogradskih okupljanja i da daje korisne informacije o njima.
Beograd nudi dosta zanimljivih sadržaja ljudima iz IT sveta.
Većina organizacija priređuje besplatna okupljanja, čini se da je jedini korak koji bi trebalo preduzeti blagovremeno informisanje o predstojećim okupljanjima.
Potrudili smo se da iznesemo dovoljno podataka kako bismo omogućili bolji uvid i odgovore na mnogobrojna pitanja koja se nameću. Neke važne podatke smo podelili, a za neke smo ostavili uputstva za pronalaženje. Svako je dobrodošao da se za dodatne informacije o ovom radu, njegovom sadržaju ili eventualnim zamerkama obrati autorima.

% Treba napomenuti da je za vreme pisanja ovog rada u republici Srbiji na snazi vanredno stanje povod pandemije virusom COVID-19. Za vreme pandemije, 
% ne održava se ni jedno okupljanje, što je trenutno posledica uredbe Vlade o okupljanjima građana za vreme navedenog vanrednog stanja.

\addcontentsline{toc}{section}{Literatura}
\bibliography{seminarski} 
\bibliographystyle{unsrt}

\end{document}
